\documentclass[10pt,a4paper]{amsart}
\usepackage{amsmath,amssymb,amsthm,graphicx}
% Retained preamble for the Stacks Project chapter {CH} (shared project preamble).
% The chapter file pinned-original-source.tex only carries \input{{preamble}} on its first line,
% so no preamble text lives inside the chapter itself.  The Stacks Project uses one single shared
% preamble for every chapter: the snapshot stages it three times (sources/stacks-chow/preamble.tex,
% sources/stacks-etale-cohomology/preamble.tex and focused-seeds/AGP/sources/stacks-snapshot-20261011/
% preamble.tex) and those copies are identical in size.  The block below is copied mechanically,
% line for line, from sources/stacks-chow/preamble.tex lines 14-18 and 154-217: the Xy-pic setup and
% the native theorem declarations together with the mathematical macros.  Nothing was added, removed
% or redefined; the wrapper class and the amsmath/amssymb/amsthm/graphicx base packages are supplied
% by the assembler.
% For commutative diagrams we use Xy-pic
% Chapter: stacks-spaces-more-morphisms.tex -- More on Morphisms of Spaces.
\usepackage[all]{xy}

% We use 2cell for 2-commutative diagrams.
\xyoption{2cell}

% Theorem environments.
%
\theoremstyle{plain}
\newtheorem{theorem}[subsection]{Theorem}
\newtheorem{proposition}[subsection]{Proposition}
\newtheorem{lemma}[subsection]{Lemma}

\theoremstyle{definition}
\newtheorem{definition}[subsection]{Definition}
\newtheorem{example}[subsection]{Example}
\newtheorem{exercise}[subsection]{Exercise}
\newtheorem{situation}[subsection]{Situation}

\theoremstyle{remark}
\newtheorem{remark}[subsection]{Remark}
\newtheorem{remarks}[subsection]{Remarks}

\numberwithin{equation}{subsection}

% Macros
%
\def\lim{\mathop{\mathrm{lim}}\nolimits}
\def\colim{\mathop{\mathrm{colim}}\nolimits}
\def\Spec{\mathop{\mathrm{Spec}}}
\def\Hom{\mathop{\mathrm{Hom}}\nolimits}
\def\Ext{\mathop{\mathrm{Ext}}\nolimits}
\def\SheafHom{\mathop{\mathcal{H}\!\mathit{om}}\nolimits}
\def\SheafExt{\mathop{\mathcal{E}\!\mathit{xt}}\nolimits}
\def\Sch{\mathit{Sch}}
\def\Mor{\mathop{\mathrm{Mor}}\nolimits}
\def\Ob{\mathop{\mathrm{Ob}}\nolimits}
\def\Sh{\mathop{\mathit{Sh}}\nolimits}
\def\NL{\mathop{N\!L}\nolimits}
\def\CH{\mathop{\mathrm{CH}}\nolimits}
\def\proetale{{pro\text{-}\acute{e}tale}}
\def\etale{{\acute{e}tale}}
\def\QCoh{\mathit{QCoh}}
\def\Ker{\mathop{\mathrm{Ker}}}
\def\Im{\mathop{\mathrm{Im}}}
\def\Coker{\mathop{\mathrm{Coker}}}
\def\Coim{\mathop{\mathrm{Coim}}}

% Boxtimes
%
\DeclareMathSymbol{\boxtimes}{\mathbin}{AMSa}{"02}

%
% Macros for moduli stacks/spaces
%
\def\QCohstack{\mathcal{QC}\!\mathit{oh}}
\def\Cohstack{\mathcal{C}\!\mathit{oh}}
\def\Spacesstack{\mathcal{S}\!\mathit{paces}}
\def\Quotfunctor{\mathrm{Quot}}
\def\Hilbfunctor{\mathrm{Hilb}}
\def\Curvesstack{\mathcal{C}\!\mathit{urves}}
\def\Polarizedstack{\mathcal{P}\!\mathit{olarized}}
\def\Complexesstack{\mathcal{C}\!\mathit{omplexes}}
% \Pic is the operator that assigns to X its picard group, usage \Pic(X)
% \Picardstack_{X/B} denotes the Picard stack of X over B
% \Picardfunctor_{X/B} denotes the Picard functor of X over B
\def\Pic{\mathop{\mathrm{Pic}}\nolimits}
\def\Picardstack{\mathcal{P}\!\mathit{ic}}
\def\Picardfunctor{\mathrm{Pic}}
\def\Deformationcategory{\mathcal{D}\!\mathit{ef}}
\newenvironment{reference}{}{}
\newenvironment{slogan}{}{}
\newenvironment{history}{}{}
\usepackage{cleveref}
\crefname{originallemma}{original lemma}{original lemmas}
\Crefname{originallemma}{Original lemma}{Original lemmas}
\makeatletter
\newlabel{perfect-lemma-pseudo-coherent-descends-fpqc}{{09UE}{source}{}{sourcepointer-1}{}}
\newlabel{perfect-lemma-pseudo-coherent-descends-fpqc@cref}{{[originallemma][0][]09UE}{[1][0][]source}{}{}{}}
\makeatother
\crefname{originallemma}{original lemma}{original lemmas}
\Crefname{originallemma}{Original lemma}{Original lemmas}
\makeatletter
\newlabel{sites-cohomology-lemma-pseudo-coherent-pullback}{{08H4}{source}{}{sourcepointer-2}{}}
\newlabel{sites-cohomology-lemma-pseudo-coherent-pullback@cref}{{[originallemma][0][]08H4}{[1][0][]source}{}{}{}}
\makeatother
\crefname{originallemma}{original lemma}{original lemmas}
\Crefname{originallemma}{Original lemma}{Original lemmas}
\makeatletter
\newlabel{spaces-perfect-lemma-derived-quasi-coherent-small-etale-site}{{071Q}{source}{}{sourcepointer-3}{}}
\newlabel{spaces-perfect-lemma-derived-quasi-coherent-small-etale-site@cref}{{[originallemma][0][]071Q}{[1][0][]source}{}{}{}}
\makeatother
\crefname{originallemma}{original lemma}{original lemmas}
\Crefname{originallemma}{Original lemma}{Original lemmas}
\makeatletter
\newlabel{spaces-perfect-lemma-descend-pseudo-coherent}{{08HE}{source}{}{sourcepointer-4}{}}
\newlabel{spaces-perfect-lemma-descend-pseudo-coherent@cref}{{[originallemma][0][]08HE}{[1][0][]source}{}{}{}}
\makeatother
\crefname{originallemma}{original lemma}{original lemmas}
\Crefname{originallemma}{Original lemma}{Original lemmas}
\makeatletter
\newlabel{spaces-topologies-lemma-recognize-fpqc-covering}{{03MS}{source}{}{sourcepointer-5}{}}
\newlabel{spaces-topologies-lemma-recognize-fpqc-covering@cref}{{[originallemma][0][]03MS}{[1][0][]source}{}{}{}}
\makeatother
\crefname{originallemma}{original lemma}{original lemmas}
\Crefname{originallemma}{Original lemma}{Original lemmas}
\makeatletter
\newlabel{spaces-topologies-lemma-refine-fpqc-schemes}{{0419}{source}{}{sourcepointer-6}{}}
\newlabel{spaces-topologies-lemma-refine-fpqc-schemes@cref}{{[originallemma][0][]0419}{[1][0][]source}{}{}{}}
\makeatother
\title{m-pseudo-coherence of a complex on a space descends along fpqc coverings}
\author{The Stacks Project Authors}
\date{stacks-project:spaces-more-morphisms.tex:local-snapshot-20261011; GNU FDL 1.2 or later}
\begin{document}
\maketitle
\noindent\textit{Editorial note.} Statement span is lines 12984-12990 of sources/stacks-spaces-more-morphisms/pinned-original-source.tex; proof span is lines 12992-13018, the complete author proof (25 body lines).\par\medskip
\noindent\textit{Editorial note.} Navigation stubs: 6 cited-but-not-included labels are supplied as non-load-bearing pointers with their pinned tags where the snapshot tag index has them.\par\medskip
\noindent\textit{Editorial note.} extra\_packages is empty on purpose: the retained preamble already loads xy with options and a second declaration clashes in the assembler wrapper.\par\medskip
\noindent\textit{Editorial note.} Text is retained verbatim, including the author's own typos and placeholders.\par\medskip
\noindent\textit{Editorial note.} Object-side (not relative) pseudo-coherence descent for spaces; pairs with the relative-base version collected on the schemes side last wave (L16096) without overlap.\par\medskip
\noindent\textit{Allowed standard prerequisites (editorial).}\begin{itemize}
\item Explicit prerequisite boundary. Schemes-side fpqc descent of pseudo-coherence (Derived Categories of Schemes), refinement of fpqc coverings of spaces by schemes (Topologies on Spaces), comparison of D\_QCoh on the small etale site with the Zariski derived category and compatibility of pseudo-coherence under it (Derived Categories of Spaces, two lemmas).
\end{itemize}\medskip
\ifdefined\hypertarget\hypertarget{sourcepointer-1}{}\fi
\noindent\textit{Source reference.} Original lemma 09UE: sources/stacks-spaces-more-morphisms/pinned-original-source.tex, source lines 13018--13018. Cited by the retained statement or proof but not part of the collected spans; supplied only as a navigation pointer. The argument does not depend on its content.\par\medskip
\ifdefined\hypertarget\hypertarget{sourcepointer-2}{}\fi
\noindent\textit{Source reference.} Original lemma 08H4: sources/stacks-spaces-more-morphisms/pinned-original-source.tex, source lines 13018--13018. Cited by the retained statement or proof but not part of the collected spans; supplied only as a navigation pointer. The argument does not depend on its content.\par\medskip
\ifdefined\hypertarget\hypertarget{sourcepointer-3}{}\fi
\noindent\textit{Source reference.} Original lemma 071Q: sources/stacks-spaces-more-morphisms/pinned-original-source.tex, source lines 13018--13018. Cited by the retained statement or proof but not part of the collected spans; supplied only as a navigation pointer. The argument does not depend on its content.\par\medskip
\ifdefined\hypertarget\hypertarget{sourcepointer-4}{}\fi
\noindent\textit{Source reference.} Original lemma 08HE: sources/stacks-spaces-more-morphisms/pinned-original-source.tex, source lines 13018--13018. Cited by the retained statement or proof but not part of the collected spans; supplied only as a navigation pointer. The argument does not depend on its content.\par\medskip
\ifdefined\hypertarget\hypertarget{sourcepointer-5}{}\fi
\noindent\textit{Source reference.} Original lemma 03MS: sources/stacks-spaces-more-morphisms/pinned-original-source.tex, source lines 13018--13018. Cited by the retained statement or proof but not part of the collected spans; supplied only as a navigation pointer. The argument does not depend on its content.\par\medskip
\ifdefined\hypertarget\hypertarget{sourcepointer-6}{}\fi
\noindent\textit{Source reference.} Original lemma 0419: sources/stacks-spaces-more-morphisms/pinned-original-source.tex, source lines 13018--13018. Cited by the retained statement or proof but not part of the collected spans; supplied only as a navigation pointer. The argument does not depend on its content.\par\medskip
\begin{lemma}
\label{lemma-pseudo-coherent-descends-fpqc}
Let $S$ be a scheme. Let $\{f_i : X_i \to X\}$ be an fpqc covering of
algebraic spaces over $S$. Let $E \in D_\QCoh(\mathcal{O}_X)$.
Let $m \in \mathbf{Z}$. Then $E$ is $m$-pseudo-coherent if and only if each
$Lf_i^*E$ is $m$-pseudo-coherent.
\end{lemma}

\begin{proof}
Pullback always preserves $m$-pseudo-coherence, see
Cohomology on Sites, Lemma
\ref{sites-cohomology-lemma-pseudo-coherent-pullback}.
Thus it suffices to assume $Lf_i^*E$ is $m$-pseudo-coherent
and to prove that $E$ is $m$-pseudo-coherent.
Then first we may assume $X_i$ is a scheme for all $i$, see
Topologies on Spaces, Lemma \ref{spaces-topologies-lemma-refine-fpqc-schemes}.
Next, choose a  surjective \'etale morphism $U \to X$ where $U$ is a scheme.
Then $U_i = U \times_X X_i$ is a scheme and we obtain an fpqc covering
$\{U_i \to U\}$ of schemes, see
Topologies on Spaces, Lemma
\ref{spaces-topologies-lemma-recognize-fpqc-covering}.
We know the result is true for
$\{U_i \to U\}_{i \in I}$ by the case for schemes, see
Derived Categories of Schemes, Lemma
\ref{perfect-lemma-pseudo-coherent-descends-fpqc}.
On the other hand, the restriction $E|_U$ comes from
an object of $D_\QCoh(\mathcal{O}_U)$ (defined using the Zariski
topology and the ``usual'' structure sheaf of $U$), see
Derived Categories of Spaces, Lemma
\ref{spaces-perfect-lemma-derived-quasi-coherent-small-etale-site}.
The lemma follows as the two notions of pseudo-coherent
(\'etale and Zariski) agree by
Derived Categories of Spaces,
Lemma \ref{spaces-perfect-lemma-descend-pseudo-coherent}.
\end{proof}

\end{document}
